%----------------------------------------------------------------------------------------
% TELLEN: SOMREGEL, PRODUCTREGEL, VARIATIES EN PERMUTATIES. Alle opgaven zijn eigen werk.
%
% Opbouw: Start, leerstof (theorie, werkwijze of uitleg, voorbeeld, korte oefening), en op een
% nieuwe pagina de Opdrachten (door elkaar, stijgende moeilijkheid), Koppel en Test jezelf.
% Labels van dit hoofdstuk beginnen met h1-.
%----------------------------------------------------------------------------------------

\mychapter{1}{Tellen}

\lkr{Leerplan B+S'': LPD 47 (telproblemen met en zonder herhaling, waarbij de volgorde al dan niet van belang is). Dit hoofdstuk behandelt de gevallen waarin de volgorde telt: herhalingsvariaties, variaties en permutaties, met de som- en productregel en het venndiagram uit de 2de graad als basis. Combinaties volgen in hoofdstuk 2. Het leerplan laat de formules toe; laat leerlingen toch eerst met de productregel redeneren, zodat de formules een samenvatting zijn en geen recept.}

Hoeveel pincodes zijn er? Op hoeveel manieren kan een klas een voorzitter en een secretaris kiezen? Hoeveel wachtwoorden moet een computer proberen om er één te kraken? Zulke vragen beantwoord je zonder alle mogelijkheden op te schrijven, door slim te tellen. Je hebt daarvoor maar twee regels nodig, de somregel en de productregel. In dit hoofdstuk tel je rijtjes waarin de volgorde belangrijk is; in hoofdstuk 2 groepjes waarin ze dat niet is.

\begin{oefeningen}[Start]
\oef{h1-start-1}

\oef{h1-verken}
\end{oefeningen}

\lkr{Laat oefening \ref{oef:h1-verken} maken vóór de uitleg. Vraag a mag met een boomdiagram, maar het patroon (vermenigvuldigen) moet eruit komen. Het verschil tussen b en c is de kern van het hoofdstuk: mag een keuze herhaald worden of niet?}

%========================================================================================
\section{De somregel en de productregel}

\subsection{De productregel}

Een outfit bestaat uit een broek, een trui en een paar schoenen. Voor elke broek kun je elke trui kiezen, en voor elk van die combinaties elk paar schoenen. Een \textbf{boomdiagram} toont dat: elke tak splitst in evenveel takken, dus vermenigvuldig je.

\begin{theorie}
\begin{stelling}{Productregel}
Bestaat een keuze uit $k$ stappen na elkaar, met $n_1$ mogelijkheden voor de eerste stap, $n_2$ voor de tweede, \ldots, $n_k$ voor de laatste, ongeacht wat er eerder gekozen is, dan zijn er

\vspace{1mm}
\qquad $\formule{n_1 \cdot n_2 \cdots n_k}$ \quad mogelijkheden.
\end{stelling}
\end{theorie}

Het woord \textit{ongeacht} is belangrijk: het aantal mogelijkheden in een stap mag niet afhangen van wat je eerder koos. De keuzes zelf mogen wel verschillen.

\voorbeeld{}{Een restaurant heeft 4 voorgerechten, 5 hoofdgerechten en 3 desserts. Hoeveel menu's met één gerecht van elk zijn er?\label{vb:h1-product}}{%
\uitleg{Een menu kies je in drie stappen. Elk voorgerecht kan met elk hoofdgerecht, en elk van die combinaties met elk dessert.}
\naast{\begin{minipage}{9.6cm}
\begin{stappenvrij}{3.4cm}{5.8cm}
4 \cdot 5 \cdot 3 = 60 & Productregel met drie stappen. Er zijn 60 menu's.
\end{stappenvrij}

\vspace{1mm}
Het boomdiagram toont het begin: elk voorgerecht splitst in 5 takken, en elk hoofdgerecht in 3.
\end{minipage}}{\begin{tikzpicture}[x=1.45cm,y=.5cm,font=\scriptsize]
  \node[blad] (s) at (0,0) {start};
  \foreach \v/\y in {1/3,2/0,3/-2,4/-4}{\node[blad] (v\v) at (1,\y) {$v_\v$}; \draw[tak] (s.east) -- (v\v.west);}
  \foreach \h/\y in {1/5.2,2/4,3/3,4/2,5/1}{\node[blad] (h\h) at (2,\y) {$h_\h$}; \draw[tak] (v1.east) -- (h\h.west);}
  \foreach \d/\y in {1/6.2,2/5.2,3/4.2}{\node[blad] (d\d) at (3,\y) {$d_\d$}; \draw[tak,thin] (h1.east) -- (d\d.west);}
  \foreach \y in {0,-2,-4}{\node[grijs] at (1.7,\y) {$\cdots$};}
  \foreach \y in {3,2,1}{\node[grijs] at (2.5,\y) {$\cdots$};}
\end{tikzpicture}}}

\subsection{De somregel}

Soms moet je eerst splitsen in gevallen die elkaar uitsluiten, zoals \textit{een code van 2 letters} of \textit{een code van 3 letters}. Dan tel je elk geval apart, en tel je op.

\begin{theorie}
\begin{stelling}{Somregel}
Valt een verzameling uiteen in delen $A$ en $B$ zonder gemeenschappelijke elementen, dan is \quad $\formule{\#(A \cup B) = \#A + \#B}$
\end{stelling}

\begin{stelling}{Somregel met overlap}
Voor eindige verzamelingen $A$ en $B$ is \quad $\formule{\#(A \cup B) = \#A + \#B - \#(A \cap B)}$
\end{stelling}

\begin{bewijs}
In $\#A + \#B$ zijn de elementen van $A \cap B$ twee keer geteld, en alle andere één keer. Trek je $\#(A \cap B)$ af, dan is elk element van $A \cup B$ precies één keer geteld.
\end{bewijs}
\end{theorie}

Met $\#A$ bedoelen we het aantal elementen van $A$. Het venndiagram toont waarom je het stuk in het midden aftrekt.

\voorbeeld{}{Hoeveel natuurlijke getallen van 1 tot en met 100 zijn deelbaar door 2 of door 5?\label{vb:h1-venn}}{%
\uitleg{Noem $A$ de veelvouden van 2 en $B$ die van 5. De getallen die in beide zitten, zijn de veelvouden van 10. We gebruiken de somregel met overlap.}
\naast{\begin{minipage}{10cm}
\begin{stappenvrij}{4.6cm}{5cm}
\#A = 50, \quad \#B = 20 & 100 : 2 = 50 en 100 : 5 = 20. \\[1ex]
\#(A \cap B) = 10 & Deelbaar door 2 én door 5: de veelvouden van 10. \\[1ex]
\#(A \cup B) = 50 + 20 - 10 = 60 & De veelvouden van 10 zijn twee keer geteld.
\end{stappenvrij}
\end{minipage}}{\begin{tikzpicture}[font=\small]
  \draw[verz] (0,0) circle (1);
  \draw[verz] (1.1,0) circle (.75);
  \node at (-.45,0) {40}; \node at (.65,0) {10}; \node at (1.45,0) {10};
  \node[above,text=accent] at (-.4,.95) {$A$}; \node[above,text=accent] at (1.5,.7) {$B$};
\end{tikzpicture}}}

Vaak is het gemakkelijker om te tellen wat je \textit{niet} wilt, en dat af te trekken van het totaal. Dat is de somregel met $A$ en het complement van $A$.

\voorbeeld{}{Hoeveel codes van 4 letters (uit 26) bevatten minstens één A?\label{vb:h1-complement}}{%
\uitleg{\textit{Minstens één A} omvat veel gevallen: één A, twee A's, \ldots{} Het tegendeel is één geval: geen enkele A. We tellen alle codes en trekken die zonder A af.}
\begin{stappen}
26^4 = 456\,976 & Alle codes: 26 mogelijkheden per plaats. \\[1ex]
25^4 = 390\,625 & Codes zonder A: 25 mogelijkheden per plaats. \\[1ex]
26^4 - 25^4 = 66\,351 & Codes met minstens één A.
\end{stappen}}

\begin{werkwijze}
Splits de keuze in stappen en vermenigvuldig (productregel). Splits in gevallen die elkaar uitsluiten en tel op (somregel); overlappen de gevallen, trek dan het gemeenschappelijke deel af. Bij \textit{minstens één} tel je vaak beter het tegendeel.
\end{werkwijze}

\nujij{Hoeveel delers heeft $72 = 2^3 \cdot 3^2$?}{Een deler is $2^a \cdot 3^b$ met $a \in \{0, 1, 2, 3\}$ en $b \in \{0, 1, 2\}$: $4 \cdot 3 = 12$ delers.}

\begin{oefeningen}
\oef{h1-product}

\oef{h1-venn}
\end{oefeningen}

%========================================================================================
\section{Rijtjes met herhaling}

Bij een pincode kies je vier keer een cijfer, en een cijfer mag opnieuw voorkomen. Elke plaats heeft dus 10 mogelijkheden. Zo'n geordend rijtje waarin herhaling toegelaten is, heet een herhalingsvariatie.

\begin{theorie}
Een \textbf{\textcolor{accent}{herhalingsvariatie}} van $p$ elementen uit $n$ is een rijtje van $p$ elementen gekozen uit $n$, waarbij de volgorde belangrijk is en elementen mogen terugkomen. Hun aantal is

\vspace{1mm}
\qquad $\formule{\HVar{n}{p} = n^p}$
\end{theorie}

Dat is de productregel met $p$ stappen van elk $n$ mogelijkheden. Je leest $\HVar{n}{p}$ als \textit{het aantal herhalingsvariaties van $p$ uit $n$}.

\voorbeeld{}{Hoeveel woorden van 5 letters bestaan er, als elk rijtje van letters meetelt, ook zonder betekenis?\label{vb:h1-hervar}}{%
\uitleg{Elke plaats krijgt een van de 26 letters, en letters mogen herhaald worden. De volgorde telt: ROMAN en MANOR zijn verschillende woorden.}
\begin{stappen}
\HVar{26}{5} = 26^5 = 11\,881\,376 & Ongeveer twaalf miljoen.
\end{stappen}}

\begin{oefeningen}
\oef{h1-hervar}
\end{oefeningen}

%========================================================================================
\section{Rijtjes zonder herhaling}

\subsection{Variaties}

Kiest een klas een voorzitter, een secretaris en een penningmeester, dan telt de volgorde: wie voorzitter is, maakt uit. Maar niemand kan twee functies krijgen. Bij elke stap is er dus één mogelijkheid minder.

\begin{theorie}
Een \textbf{\textcolor{accent}{variatie}} van $p$ elementen uit $n$ (met $p \leqslant n$) is een rijtje van $p$ verschillende elementen gekozen uit $n$, waarbij de volgorde belangrijk is. Hun aantal is

\vspace{1mm}
\qquad $\formule{\Var{n}{p} = n(n - 1)(n - 2) \cdots (n - p + 1)}$ \qquad ($p$ factoren)
\end{theorie}

\voorbeeld{}{Een klas van 22 leerlingen kiest een voorzitter, een secretaris en een penningmeester. Op hoeveel manieren kan dat?\label{vb:h1-var}}{%
\uitleg{De volgorde telt, want de functies zijn verschillend, en niemand krijgt twee functies. Dat is een variatie van 3 uit 22.}
\begin{stappen}
\Var{22}{3} = 22 \cdot 21 \cdot 20 = 9\,240 & 22 keuzes voor de voorzitter, dan 21 voor de secretaris, dan 20 voor de penningmeester.
\end{stappen}}

\subsection{Permutaties}

Kies je alle $n$ elementen, dan rangschik je ze gewoon. Het product $n(n - 1) \cdots 2 \cdot 1$ komt zo vaak voor dat het een eigen naam heeft.

\begin{theorie}
De \textbf{\textcolor{accent}{faculteit}} van een natuurlijk getal $n \geqslant 1$ is \quad $\formule{n! = n(n - 1) \cdots 2 \cdot 1}$, \quad en we spreken af dat $\formule{0! = 1}$.

\vspace{1mm}
Een \textbf{\textcolor{accent}{permutatie}} van $n$ elementen is een rangschikking van alle $n$ elementen. Hun aantal is \quad $\formule{\Perm{n} = n!}$

\begin{stelling}{Variaties met faculteiten}
Voor $0 \leqslant p \leqslant n$ is \quad $\formule{\Var{n}{p} = \dfrac{n!}{(n - p)!}}$
\end{stelling}

\begin{bewijs}
In $\frac{n!}{(n - p)!}$ vallen de factoren $(n - p)(n - p - 1) \cdots 1$ weg tegen de noemer, en blijven $n(n - 1) \cdots (n - p + 1)$ over: precies $\Var{n}{p}$.
\end{bewijs}
\end{theorie}

Je leest $n!$ als \textit{$n$ faculteit}. Faculteiten groeien heel snel: $10! = 3\,628\,800$ en $20!$ heeft al 19 cijfers. De afspraak $0! = 1$ zorgt ervoor dat de formule ook klopt voor $p = n$: $\Var{n}{n} = \frac{n!}{0!} = n!$.

\voorbeeld{}{Zes vrienden gaan naast elkaar zitten in de bioscoop. Op hoeveel manieren kan dat? En als Anna en Ben naast elkaar willen zitten?\label{vb:h1-perm}}{%
\uitleg{Alle zes worden gerangschikt: een permutatie. Voor de tweede vraag behandelen we Anna en Ben eerst als één blok, en tellen daarna op hoeveel manieren ze binnen het blok zitten.}
\begin{stappen}
\Perm{6} = 6! = 720 & Alle rangschikkingen. \\[1ex]
\Perm{5} \cdot \Perm{2} = 120 \cdot 2 = 240 & Het blok \textit{Anna-Ben} en de vier anderen: 5 eenheden. Binnen het blok: Anna links of Ben links.
\end{stappen}}

\begin{werkwijze}
Vraag je af: telt de volgorde, en mag een element terugkomen? Volgorde en herhaling: $n^p$. Volgorde zonder herhaling: $\Var{n}{p}$, of $n!$ als je alle elementen rangschikt. Bij een voorwaarde zoals \textit{naast elkaar}: neem de elementen samen als één blok, en vermenigvuldig met het aantal rangschikkingen binnen het blok.
\end{werkwijze}

\nujij{Bereken $\Var{7}{2}$ en $\dfrac{9!}{7!}$.}{$\Var{7}{2} = 7 \cdot 6 = 42$ en $\frac{9!}{7!} = 9 \cdot 8 = 72$.}

\begin{oefeningen}
\oef{h1-var}
\end{oefeningen}

%========================================================================================
\opdrachtenkop

\begin{oefeningen}
\oef{h1-o1}

\oef{h1-o2}

\oef{h1-o3}

\oef{h1-o4}

\oef{h1-o5}

\oef{h1-o6}

\oef{h1-o7}

\oef{h1-o8}

\oef{h1-o9}

\oef{h1-o10}

\oef{h1-o11}

\oef{h1-o12}

\oef{h1-o13}

\oef{h1-o14}
\end{oefeningen}

%----------------------------------------------------------------------------------------
% KOPPEL: verbanden met andere leerstof
%----------------------------------------------------------------------------------------
\needspace{20\baselineskip}
\opdrachtsectie{Koppel}

\lkr{De eerste oefening gebruikt de wet van Laplace (kans = gunstig over mogelijk), uit de 2de graad; Kansrekenen staat in hetzelfde jaar. De tweede gebruikt logaritmen uit de cursus Exponentiële en logaritmische functies (5de jaar).}

\begin{oefeningen}
\oef{h1-koppel-1}

\oef{h1-koppel-2}
\end{oefeningen}

%----------------------------------------------------------------------------------------
\opdrachtsectie{Test jezelf}

\begin{oefeningen}
\oef{h1-t1}

\oef{h1-t2}

\oef{h1-t3}

\oef{h1-t4}

\oef{h1-t5}

\oef{h1-t6}
\end{oefeningen}

\herhaalschema{%
  \fv{oef:h1-t1}{vb:h1-product}{oef:h1-product}%
  \fv{oef:h1-t2}{vb:h1-venn}{oef:h1-venn}%
  \fv{oef:h1-t3}{vb:h1-hervar}{oef:h1-hervar}%
  \fv{oef:h1-t4}{vb:h1-var}{oef:h1-var}%
  \fv{oef:h1-t5}{vb:h1-perm}{oef:h1-var}%
  \fv{oef:h1-t6}{vb:h1-complement}{oef:h1-venn}}

\ifzelfstudie
\section*{Maak je eigen samenvatting}
Sluit dit boek en schrijf in je schrift: de productregel en de somregel (ook met overlap), wanneer je het tegendeel telt, wat herhalingsvariaties, variaties en permutaties zijn met hun formules, wat $n!$ is en waarom $0! = 1$, en hoe je een voorwaarde als \textit{naast elkaar} aanpakt. Vergelijk daarna met de kaders.
\fi
